\documentclass[11pt]{article}
\usepackage[margin=1in]{geometry}
\usepackage{setspace}
\usepackage{parskip}
\usepackage{booktabs}
\usepackage{array}
\usepackage{tabularx}
\usepackage{caption}
\usepackage{titlesec}
\usepackage{fancyhdr}
\usepackage{hyperref}
\usepackage{amsmath}
\usepackage{microtype}
\usepackage{enumitem}
\usepackage{xcolor}
\usepackage{graphicx}
\definecolor{reportblue}{RGB}{46,64,87}
\hypersetup{
    colorlinks=true,
    linkcolor=reportblue,
    urlcolor=reportblue,
    citecolor=reportblue
}
\titleformat{\section}{\large\bfseries\color{reportblue}}{}{0em}{}[\titlerule]
\titleformat{\subsection}{\normalsize\bfseries\color{reportblue}}{}{0em}{}
\titleformat{\subsubsection}{\normalsize\bfseries}{}{0em}{}
\titlespacing*{\section}{0pt}{18pt}{8pt}
\titlespacing*{\subsection}{0pt}{12pt}{4pt}
\titlespacing*{\subsubsection}{0pt}{10pt}{4pt}
\pagestyle{plain}
\fancyhf{}
\rfoot{\small\thepage}
\renewcommand{\headrulewidth}{0.4pt}
\captionsetup[table]{position=above, labelfont=bf, labelsep=period}
\captionsetup[figure]{position=below, labelfont=bf, labelsep=period}
\setlength{\parindent}{0pt}
\setlength{\parskip}{8pt}
% ── Column types for tables ───────────────────────────────────────────────────
\newcolumntype{L}[1]{>{\raggedright\arraybackslash}p{#1}}

\setstretch{1.15}
\setlist[itemize]{noitemsep, topsep=4pt, left=1.5em}

\begin{document}
\pagenumbering{gobble}

% ── TITLE PAGE ────────────────────────────────────────────────────────────────
\begin{titlepage}
  \centering
  \vspace*{3cm}
  {\LARGE\bfseries\color{reportblue} Big Belly Smart Bin Forecasting\par}
  \vspace{0.5cm}
  {\large Cal Zero Waste - UC Berkeley Facilities Services\par}
  \vspace{2.5cm}
  {\normalsize Joy Zhang \quad Jacob Quisumbing \quad Vicky Xiao \quad Zoey Kim\par}
  \vspace{0.4cm}
  {\normalsize IND ENG 180 - Project 10A\par}
  \vspace{0.4cm}
  {\normalsize University of California, Berkeley\par}
  \vspace{0.4cm}
  {\normalsize April 10, 2026\par}
\end{titlepage}

\pagenumbering{arabic}
\setcounter{page}{1}

% ── EXECUTIVE SUMMARY ─────────────────────────────────────────────────────────
\section*{Executive Summary}
\addcontentsline{toc}{section}{Executive Summary}

Cal Zero Waste, a unit within UC Berkeley's Facilities Services and Campus Operations
department, manages 75 Big Belly smart bin stations - approximately 248 individual bins -
deployed across the main Berkeley campus. Despite sensor infrastructure that reports each
bin's fill level in real time, collection operations remain largely reactive: trucks are
dispatched only after a bin triggers an alert, or after seven days have passed regardless
of how full the bin is.

Analysis of 59,111 collection events from January 2023 through February 2026 shows that
35.1\% of all collections were unnecessary - bins were emptied well before they needed to
be. On average, bins were only 47.1\% full at the time of collection, despite a 60\%
trigger threshold. The division's collection efficiency metric has averaged approximately
61\% over the past six months, reflecting this systematic over-collection.

This project builds a forecasting tool that predicts when each bin will next need
collection - up to one week in advance. Rather than waiting for a bin to trigger an alert,
Cal Zero Waste would know in advance which bins are approaching their threshold and which
can safely wait. The tool also classifies each bin by how frequently it needs service:
roughly one quarter of bins need more than one collection per week, most bins need
approximately one collection per week, and a small number rarely fill up under normal
conditions.

The forecasting approach performs well for weekly scheduling - it correctly identifies 95\%
of bins that will need collection within the next seven days. It is less reliable for
predicting which specific bins will need service the following day, a limitation addressed
by a complementary short-term fill estimation model. Together, the two tools are designed
to supplement the existing CLEAN dashboard workflow, not replace it.

The uncertainty bounds accompanying each forecast have been empirically validated and
calibrated to around 80\% coverage using a held-out calibration set. Further work in progress
includes coordination with Project 10B on how these predictions feed into dispatch routing.

\newpage

% ── INTRODUCTION ──────────────────────────────────────────────────────────────
\section{Introduction}

\subsection{Organization Overview}

Cal Zero Waste is the waste management program of UC Berkeley's Facilities Services and
Campus Operations department, responsible for refuse, compost, and recycling collection
across the main Berkeley campus. Cal Zero Waste employs a small operations team that
coordinates daily collection routes and manages the Big Belly smart bin network. The
project sponsor and primary point of contact is Samantha Bunke, Zero Waste Specialist.

\subsection{The Big Belly System}

Each Big Belly station is a solar-powered kiosk with a compacting capacity of approximately
150 gallons. Fill level refers to the percentage of that compacted capacity that has been
used - a bin at 60\% fill level contains the equivalent of 90 gallons of compacted waste.
Sensors inside each bin continuously measure fill level and transmit that data to Big
Belly's CLEAN Management Console, a cloud-based dashboard that shows the current status of
every bin across the network.

Each station contains three separate bins - one for landfill waste, one for compostables,
and one for recyclables - for a total of approximately 248 individual bins across 90 campus
locations. Each waste stream is collected separately, often on different schedules, by
different trucks. A single station can therefore have a landfill bin that needs collection
urgently while its adjacent compost bin is nearly empty.

Collection is currently triggered by one of two conditions: a bin reaches 60\% of its
compacted capacity, or seven days have passed since the last collection, whichever occurs
first. Recyclables bins are an exception - their age-based trigger is disabled, and
collection is triggered by fullness only. When a bin reaches threshold, the CLEAN dashboard
displays a red alert. Drivers follow predefined routes that are adjusted on the fly in
response to these alerts, visiting a bin as part of a route whether or not it currently
needs service.

\subsection{Motivation}

The alert-based system produces two types of operational failure. Over-collection occurs
when bins are serviced before they need to be: either because a driver is nearby and
empties the bin opportunistically, or because a route includes the bin on a fixed schedule
regardless of its fill level. Under-collection occurs when bins in high-traffic areas reach
capacity between scheduled route visits, producing overflow incidents. Analysis of the full
dataset confirms both problems at scale: 35.1\% of collections are wasted trips, and
average bin fullness at the time of collection is only 47.1\%. This project develops
forecasting models to give Cal Zero Waste advance visibility into bin fill levels, enabling
a shift toward proactive scheduling.

\newpage

% ── PROBLEM STATEMENT ─────────────────────────────────────────────────────────
\section{Problem Statement}

\subsection{Problem Description}

The core problem is that Cal Zero Waste has no way to predict future bin fill levels. The
CLEAN dashboard shows what each bin's fill level is right now, but not what it will be
tomorrow or next week. Without that forward-looking information, operators cannot
distinguish between a bin that will reach threshold in two hours and one that will not
reach it for five days. This forces either over-servicing - visiting bins before they need
it - or under-servicing - waiting for alerts that may arrive too late.

The problem is compounded by variation across the bin network. Fill rates vary by a factor
of five to ten across campus locations. Landfill bins fill most quickly and are most
sensitive to foot traffic and event surges. Compostable bins follow dining schedule
patterns. Recyclables bins fill most slowly and operate under different service rules. A
single station can simultaneously have a full landfill bin and a nearly empty compost bin,
meaning any useful forecast must be generated at the individual bin level, not the station
level.

\subsection{What We Are Predicting}

For each bin in the network and each waste stream it serves, the goal is to predict how
full the bin will be at future points in time, and specifically, how many hours remain
until it reaches the 60\% collection threshold. This time-to-threshold estimate - referred
to as $T^*$ - is the primary operational output. It tells a scheduler: ``this bin will
need service in approximately $X$ hours.'' The formal definition is:
\[
  T^*_{i,k} = \min\bigl\{t : \hat{F}_{i,k}(t) \geq 0.60\bigr\}
\]
where $\hat{F}_{i,k}(t)$ is the predicted fill level of bin $i$ of stream type $k$ at
time $t$. Each prediction also carries an uncertainty range - an optimistic estimate (the
bin could fill sooner) and a conservative estimate (the bin could fill later) - so that
schedulers can make informed decisions under uncertainty.

\subsection{Assumptions and Constraints}

The primary dataset covers January 2023 through February 2026, containing 59,111
collection events across 248 bins. Data from 2020 and 2021 is excluded due to
pandemic-related disruption to campus occupancy and waste generation patterns, which
produced fill dynamics unrepresentative of normal operations. Each bin is assumed to have
a consistent 150-gallon compacting capacity, and the current bin network configuration is
assumed stable over the forecasting horizon. Fill level is treated as resetting to
approximately 0\% after each collection event. The 60\% fullness threshold and seven-day
age rule are treated as fixed operational parameters.

The main data constraint is that the CLEAN system can only export collection event logs -
not a continuous record of fill level between collections. This means fill levels can only
be observed at the moment a bin is emptied, and even then only in coarse steps of 0, 20,
40, 60, 80, or 100\%. This limits how precisely the model can characterize fill dynamics
and introduces noise that cannot be removed through modeling alone.

\subsection{Success Criteria}

Because missing a bin that needs collection carries greater operational risk than
unnecessarily visiting one that does not, the forecasting model is designed to favor
catching all bins that need service, even at the cost of some false alarms. The division's
stated goal is to improve collection efficiency from the current 60-70\% range to
approximately 90\%. Achieving this depends on both the accuracy of the forecast (this
project) and how the forecast is used to plan collection routes (Project 10B).

\subsection{External Factors}

Several factors drive the variation in fill rates that makes forecasting valuable. The UC
Berkeley academic calendar is the largest driver: campus population drops sharply during
breaks and summer and rises sharply during move-in week and the start of each semester.
High-impact events - Cal Day, the April 20 gathering at Memorial Glade, commencement, and
home football games - produce localized waste spikes at nearby bins. Weather affects how
many people are outside and therefore how quickly outdoor bins fill. Operational changes,
such as the reusable container program introduced at MLK Student Union, can also shift
fill-rate patterns for affected bins.

\newpage

% ── PROBLEM ANALYSIS AND SOLUTION ────────────────────────────────────────────
\section{Problem Analysis and Solution}

\subsection{Data Collection}

The primary dataset was exported from the CLEAN Management Console by the project sponsor.
Three daily collection activity logs cover January 2023 through February 2026, containing
59,111 events across 248 bins at 90 campus locations. Each record includes bin serial
number, location, waste stream type, reason for collection, fill level at collection (in
20\% increments), and a timestamp.

Supplementary data includes UC Berkeley academic calendars for 2023-24, 2024-25, and
2025-26 from the Office of the Registrar; Cal Bears football home game schedules for
2023-2025; and daily weather data for Berkeley from the Open-Meteo historical archive API.
The team also interviewed Charles Montgomery, Park Maintenance Supervisor for the City of
Walnut Creek, whose department manages 35 Big Belly stations at a higher collection
threshold (75-80\%) with per-bin adjustments based on material weight.

\subsection{Exploratory Data Analysis}

Analysis of the full collection log confirms the scale of operational inefficiency driving
this project. Of the 59,111 collection events, 35.1\% are classified as Not Ready -
meaning the bin was collected well below the 60\% threshold. Average fill level at the
time of collection is 47.1\%, meaningfully below the operational target.
Table~\ref{tab:reasons} breaks down collection events by reason.

\begin{table}[ht]
\centering
\caption{Collection event breakdown by reason, January 2023 through February 2026.}
\label{tab:reasons}
\begin{tabular}{lrr}
\toprule
\textbf{Collection Reason}   & \textbf{Event Count} & \textbf{Share of Total} \\
\midrule
Fullness                     & 31,888               & 53.9\%                  \\
Not Ready (wasted trip)      & 20,742               & 35.1\%                  \\
Age (7-day rule)             & 6,481                & 11.0\%                  \\
\midrule
\textbf{Total}               & \textbf{59,111}      & \textbf{100.0\%}        \\
\bottomrule
\end{tabular}
\end{table}

\subsubsection*{Wasted Trip Analysis}

To better understand which wasted trips are actually preventable, the team categorized
each of the 20,742 Not Ready collections into one of four groups:

\begin{itemize}
  \item \textbf{Pre-event staging} (4,761 trips, 23.0\% of wasted trips): Collections
  that occurred within six days before a known high-impact campus event or semester start.
  These may reflect intentional emptying before an anticipated surge. The six-day window
  was selected empirically by testing windows of one to seven days and choosing the window
  that produced the largest lift in Not Ready rate relative to the baseline. These trips
  are only partially eliminable - Cal Zero Waste may deliberately empty bins before Cal
  Day regardless of fill level, which is operationally sensible.

  \item \textbf{Frequency mismatch} (3,736 trips, 18.0\% of wasted trips): Collections
  at bins whose historical fill rate indicates they do not need more than one collection
  per week, but that were visited more frequently. These trips are directly eliminable
  through better scheduling.

  \item \textbf{Unexplained} (12,245 trips, 59.0\% of wasted trips): Collections that
  cannot be attributed to event staging or frequency mismatch. These represent the largest
  category and are the most directly addressable by the forecasting model.

  \item \textbf{Age-rule forced} (0 trips): No wasted trips were attributable solely to
  the seven-day age rule, indicating that bins collected on the age rule were generally
  full enough to justify service.
\end{itemize}

Table~\ref{tab:wasted} summarizes the breakdown. Approximately 77\% of wasted trips are
potentially eliminable through better forecasting and scheduling, representing roughly
27\% of all collections. The quarterly trend is also notable: the unexplained wasted trip
rate has risen from approximately 8\% of all collections in 2023 Q3 to approximately 29\%
in early 2026, suggesting that the gap between current scheduling practices and actual bin
needs has widened over time.

\begin{table}[ht]
\centering
\caption{Wasted trip breakdown by category, January 2023 through February 2026.}
\label{tab:wasted}
\begin{tabularx}{\textwidth}{L{3.2cm} r r r L{3.8cm}}
\toprule
\textbf{Category} & \textbf{Count} & \textbf{\% of Wasted} & \textbf{\% of All} &
\textbf{Eliminable by Forecast} \\
\midrule
Pre-event staging   & 4,761  & 23.0\% & 8.1\%  & Partially (depends on intent) \\
Frequency mismatch  & 3,736  & 18.0\% & 6.3\%  & Yes \\
Unexplained         & 12,245 & 59.0\% & 20.7\% & Yes \\
Age-rule forced     & 0      & 0.0\%  & 0.0\%  & No (requires policy change) \\
\midrule
\textbf{Total}      & \textbf{20,742} & \textbf{100\%} & \textbf{35.1\%} & \\
\bottomrule
\end{tabularx}
\end{table}

Fill rates also vary substantially by stream type and location. Landfill bins fill most
quickly and are most event-sensitive. Compostables bins follow dining patterns.
Recyclables bins fill most slowly, with a higher proportion rarely reaching threshold
under typical conditions. Day-of-week patterns are pronounced: needed collection counts
drop sharply on weekends, consistent with lower campus occupancy.

\subsection{Methodology}

The modeling work proceeds in three phases. The first fits curves to each bin's historical
fill data to characterize how that bin typically fills over time. The second trains
statistical models to predict when a given bin will next need collection. The third
combines both outputs into a rolling daily forecast that gives the operations team a
forward-looking view of the entire bin network.

\subsubsection*{Phase 1: Fill-Curve Modeling}

Because the CLEAN system does not record fill levels between collection events, the team
infers each bin's fill trajectory from the relationship between how long has passed since
the last collection and how full the bin was when it was next collected. Five curve shapes
were tested for each bin and the best-fitting shape was selected. The S-shaped curve
(sigmoid) was selected for 82\% of bins. The plateau at high fill levels likely reflects a behavioral effect. As bins approach capacity, users dispose of waste elsewhere rather than forcing material in rather than compaction mechanics alone. The sensor reports fill level as a percentage of compacting capacity, so the plateau may also reflect a measurement artifact: compaction continues at high levels but the reported percentage changes more slowly as the compactor works against denser material. Both mechanisms produce the same observable sigmoid shape in the data and are not distinguishable from collection-event-level observations alone.
For the 14 low-usage bins with too few observations to fit a reliable curve, the model
uses the median fill rate for that bin's waste stream as a fallback. The average
goodness-of-fit ($R^2 = 0.23$) reflects the coarse 20\% measurement increments in the
data rather than a failure of the curves - individual collection events vary substantially
around the average trajectory, and this variability cannot be reduced through modeling
alone.

\subsubsection*{Phase 2: Statistical Models}

Four models were trained and evaluated. All models used 2023-2024 data for training and
2025 onward for testing. This is a temporal split that prevents the model from being evaluated
on data it was trained on. The quantity to predict for the regression models is the number of
hours until the next collection event for a given bin. For the binary classification
model, the two classes are: class 1 (the bin's fill level is currently at or above the
60\% collection threshold and needs service); class 0 (the bin is below threshold and does
not yet need service).

The inputs used to make predictions include: how much time has passed since the bin was
last collected; how full it was at that last collection; the day of the week and time of
year; whether the campus is currently in session; whether a high-impact event is
approaching; weather conditions; and each bin's historical fill characteristics derived
from Phase 1.

\begin{itemize}
  \item The \textbf{baseline model} predicts that each bin will take the same amount of
  time to need collection as it historically has on that day of the week. This is the
  minimum standard all other models must exceed.

  \item The \textbf{daily demand forecast} (Prophet) predicts how many total collections
  each waste stream will need on a given day across the entire network. This is useful for
  staffing and truck scheduling at the fleet level but does not indicate which specific
  bins need service.

  \item The \textbf{binary readiness classifier} (XGBoost) predicts whether a given bin has reached the collection threshold at the time of its most recent recorded event. It is tuned to catch as many bins that need collection as
  possible, accepting some false positives to minimize missed bins.

  \item The \textbf{quantile forecast} (LightGBM) predicts three estimates of when a bin
  will need collection: an optimistic case, a median estimate, and a conservative case. These bounds give the operations team a range to work with rather than a single number.
\end{itemize}


Training data for all statistical models is restricted to true fill cycles. Rows where
the subsequent collection reason is Not Ready are excluded, so the models learn to predict
the interval to the next needed collection rather than to the next scheduled visit,
regardless of whether that visit was warranted.

\subsection{Model Interpretation for Operations Use}
The forecasting system is designed to be interpretable and usable by operations staff without a technical background. At a high level, the system answers two core questions: how full each bin is likely to be right now, and when it is expected to require collection. Rather than relying on a single model, the system combines multiple components that each serve a specific purpose in decision-making.

The fill-curve model captures how each bin typically fills over time. Because the available data only records fill level at the moment of collection, the model learns the average relationship between time since last service and observed fullness. This allows the system to estimate a bin’s current fill level even when no recent sensor reading is available. In practice, this functions as a “baseline expectation” for how quickly a bin fills under normal conditions and is particularly useful for short-term urgency assessment.

The machine learning forecast builds on this by predicting when a bin will next reach the 60\% threshold. It incorporates recent service history, day-of-week patterns, seasonal trends, and external factors such as campus events. Rather than assuming each bin behaves the same over time, the model adjusts predictions dynamically based on recent observations, making it more responsive to changes in usage patterns. This component is most valuable for planning several days to one week in advance.

The binary readiness classifier simplifies the decision problem into a yes-or-no question: does this bin need collection now? It is intentionally tuned to prioritize recall, meaning it is designed to catch as many bins that require service as possible, even if that results in some unnecessary flags. This reflects the operational reality that missing a full bin is more costly than servicing one slightly early.

Finally, the daily demand forecast operates at the system level by estimating how many bins will require collection across the network on a given day. While it does not identify specific bins, it provides important context for staffing and route planning by indicating expected workload levels.

Together, these components are translated into a simplified urgency labeling system that maps technical predictions into actionable categories. This design ensures that the forecasting system supports decision-making without requiring operators to interpret model outputs directly.

\subsection{Results}

\subsubsection*{How Accurately Can We Predict When a Bin Will Need Collection?}

The quantile forecast model reduces prediction error by 16\% relative to the baseline
(mean error of 40.3 hours versus 48.3 hours). To put this in context: the average bin
needs collection roughly every two to three days, so an average error of 40 hours
represents roughly half a collection cycle. The model is more accurate for once-a-week
bins with stable fill patterns and less accurate for high-traffic bins with variable fill
rates.

\begin{table}[ht]
\centering
\caption{Prediction accuracy on the 2025 test set. The target is the number of hours
until the bin next needs collection.}
\label{tab:regression}
\begin{tabular}{lrr}
\toprule
\textbf{Model}    & \textbf{Mean Error (hours)} & \textbf{Typical Error (hours)} \\
\midrule
Baseline          & 48.3                        & 65.8                           \\
Quantile Forecast & 40.3                        & 54.3                           \\
\bottomrule
\end{tabular}
\end{table}

% Weather features contributed minimally to prediction accuracy, reflecting Berkeley's mild
% climate. The strongest predictive signals are the bin's own historical fill
% characteristics, how much time has passed since it was last collected, and the day of the
% week.

\subsubsection*{Feature Importance and Per-Bin Error}

\begin{figure}[ht]
    \centering
    \includegraphics[width=0.85\textwidth]{feature_importance_q50.png}
    \caption{Top 20 gain-based features for the LightGBM Q50 model.}
    \label{fig:feature_importance_q50}
\end{figure}

Figure 1 shows the gain-based feature importance for the LightGBM Q50 model. The model relies most heavily on recency and recent fill-history variables. The three highest-gain predictors are the 28-day rolling average of elapsed time, the 7-day rolling average of elapsed time, and the current elapsed hours since the bin was last collected. Bin-specific historical fill characteristics also rank prominently, including median elapsed hours, rolling 7-day fullness, projected fill at 24, 48, and 72 hours, mean fullness at collection, and location. Weather variables appear among the top gain-based features, particularly wind and daily temperature, but the overall pattern remains consistent with the central result of this project: the most useful predictive information is a bin’s own recent service history and historical fill behavior.

\begin{figure}[ht]
    \centering
    \includegraphics[width=0.75\textwidth]{per_bin_error_histogram.png}
    \caption{Distribution of per-bin prediction error (mean absolute error in hours) for the LightGBM Q50 model.}
    \label{fig:per_bin_error_hist}
\end{figure}


Aggregate accuracy metrics also mask substantial heterogeneity across the fleet. Across the 206 bins with sufficient test observations for bin-level evaluation, the median per-bin mean absolute error is 34.9 hours, with an interquartile range of 28.4 to 43.3 hours and a full range of 15.5 to 92.3 hours. A total of 136 of 206 bins have per-bin mean absolute error at or below 40 hours, while only 12 exceed 60 hours. Figure 2 shows that most bins cluster in the 25-45 hour range, with a smaller right tail of bins that are materially harder to predict.

\begin{figure}[ht]
    \centering
    \includegraphics[width=0.85\textwidth]{hardest_bins_q50.png}
    \caption{Fifteen hardest bins to predict under the LightGBM Q50 model.}
    \label{fig:hardest_bins_q50}
\end{figure}

The hardest-to-predict bins are concentrated in the recyclables stream. The 15 worst bins include E05 V\&A Cafe/Etcheverry, 001 Berkeley Way (Offsite), E13 Stanley Hall, E14 Physics, and E12 Evans Hall Bus Stop; most are recyclables bins with long and irregular collection cycles. By stream, mean per-bin mean absolute error is 49.8 hours for Recyclables, compared with 34.0 hours for Compostables and 30.1 hours for Landfill. This reinforces the broader interpretation of Table 3: the model performs best on bins with regular service rhythms and is less precise on bins with sparse or more variable demand over time.

\subsubsection*{How Well Can We Predict Total Daily Collection Demand?}

The daily demand forecast predicts how many bins across the network will need service each
day by waste stream. Table~\ref{tab:prophet} shows the average prediction error on the
2025 test period.

\begin{table}[ht]
\centering
\caption{Daily demand forecast accuracy by waste stream, 2025 test period. Values are
collections per day.}
\label{tab:prophet}
\begin{tabular}{lrr}
\toprule
\textbf{Stream}  & \textbf{Mean Error} & \textbf{Typical Error} \\
\midrule
Landfill         & 7.7                 & 10.2                   \\
Compostables     & 6.8                 & 8.6                    \\
Recyclables      & 6.9                 & 9.1                    \\
\bottomrule
\end{tabular}
\end{table}


\subsubsection*{How Well Can We Identify Which Bins Need Collection?}

The binary readiness classifier predicts whether a given bin is currently at or above the
60\% collection threshold. At the default operating point, the model correctly identifies
74\% of bins that need collection. By adjusting the operating point to be more aggressive,
that figure improves to 85\% - at the cost of also flagging some bins that do not yet need
service. Table~\ref{tab:xgboost} summarizes both operating points.

\begin{table}[ht]
\centering
\caption{Binary readiness classifier performance, 2025 test set. Recall is the fraction
of bins that actually need collection that the model correctly identifies.}
\label{tab:xgboost}
\begin{tabular}{L{4.5cm}cccc}
\toprule
\textbf{Operating Point} & \textbf{AUC} & \textbf{Precision} & \textbf{Recall} &
\textbf{F1} \\
\midrule
Default (0.50)           & 0.857        & 0.77               & 0.74            & 0.75 \\
Recall-optimized (0.33)  & 0.857        & 0.70               & 0.85            & 0.77 \\
\bottomrule
\end{tabular}
\end{table}

Given the emphasis on minimizing missed collections, the recall-optimized operating point
is recommended for operational use.

\subsubsection*{How Far Ahead Can the Model Reliably Predict?}

Table~\ref{tab:horizons} shows how well the quantile forecast identifies bins that will
need collection within various time windows.

\begin{table}[ht]
\centering
\caption{Forecast reliability by planning horizon, 2025 test set. Recall is the fraction
of bins that will actually need collection within the horizon that the model correctly
flags.}
\label{tab:horizons}
\begin{tabular}{L{2.5cm}ccc}
\toprule
\textbf{Horizon} & \textbf{AUC} & \textbf{Precision} & \textbf{Recall} \\
\midrule
24 hours         & 0.737        & 0.702              & 0.122           \\
48 hours         & 0.771        & 0.815              & 0.371           \\
72 hours         & 0.809        & 0.870              & 0.579           \\
1 week           & 0.871        & 0.904              & 0.948           \\
\bottomrule
\end{tabular}
\end{table}

At one week out, the model correctly identifies 95\% of bins that will need service
within the next seven days. This is a reliable basis for weekly scheduling decisions. At 24
hours, recall drops to 12\%, meaning the model misses approximately 88\% of bins needing
collection the following day. This is a structural limitation of the data: because fill
levels are only recorded when a bin is collected, the model has no way to observe what
happens between service visits. The fill-curve model from Phase 1 is better suited for
the 0-48 hour window because it uses the bin's characteristic fill trajectory to estimate
its current state directly.

The two approaches are complementary: the statistical model is most useful for planning
the week ahead and classifying which bins need more or less frequent service; the
fill-curve model is most useful for near-term urgency assessment.

\subsection{Operational Forecast Output}

The modeling outputs are combined into a rolling daily forecast that gives Cal Zero Waste
a forward-looking view of every bin in the network. For each of the 248 bins, the forecast
provides:

\begin{itemize}
  \item An estimated current fill level, based on how long has passed since the last
  collection and that bin's typical fill trajectory
  \item A projected date and time when the bin is expected to reach the 60\% threshold,
  along with an earlier and later bound reflecting uncertainty
  \item An urgency label combining the fill-curve estimate and the optimistic bound from
  the quantile model
  \item A service frequency classification indicating whether the bin typically needs
  more than one collection per week, approximately one per week, or rarely needs service
\end{itemize}

Table~\ref{tab:urgency} defines the rules mapping model outputs to each urgency label.
The conditions are evaluated in order from top to bottom; the first matching condition
determines the label.

\begin{table}[ht]
\centering
\caption{Urgency label definitions. Conditions are evaluated in priority order.
``Days until full (curve)'' is the fill-curve point estimate; ``Days until full (optimistic)''
is the P10 bound from the quantile model.}
\label{tab:urgency}
\begin{tabularx}{\textwidth}{L{2cm} X}
\toprule
\textbf{Label} & \textbf{Condition} \\
\midrule
OVERDUE  & Current fill level $\geq$ 60\% (bin is already at or past threshold) \\
MONITOR  & No curve fit available, or fill curve asymptote is below 60\% (bin rarely
           reaches threshold; check periodically) \\
HIGH     & Days until full (optimistic) $\leq$ 1 (even the fast-fill scenario says
           collect within 24 hours) \\
ELEVATED & Days until full (curve) $\leq$ 3 (expected to reach threshold within 3 days) \\
MODERATE & Days until full (curve) $\leq$ 5 (expected to reach threshold within 5 days) \\
LOW      & Days until full (curve) $>$ 5 (no collection needed within the current week) \\
\bottomrule
\end{tabularx}
\end{table}

Table~\ref{tab:frequency} shows the service frequency breakdown across the network.

\begin{table}[ht]
\centering
\caption{Service frequency classification across the 248-bin network, based on fill-curve
estimates of time to reach the 60\% threshold from empty.}
\label{tab:frequency}
\begin{tabularx}{\textwidth}{L{2.8cm} r r X}
\toprule
\textbf{Class} & \textbf{Count} & \textbf{Share} & \textbf{Description} \\
\midrule
Multi-visit  & 63  & 25\% & Reaches 60\% in 3.5 days or less; needs more than one
                            collection per week \\
Once-a-week  & 147 & 59\% & Reaches 60\% in 3.5-7 days; one weekly collection is
                            sufficient \\
Unmatched    & 14  & 6\%  & In CLEAN log but absent from fill-curve dataset;
                            classified as MONITOR until sufficient history is available \\
\bottomrule
\end{tabularx}
\end{table}

As of the most recent data timestamp (February 20, 2026), the forecast flags 80 bins as
requiring collection immediately, and a total of 142 bins as needing collection at some
point during the current week (inclusive of the 80 flagged immediately). Median time to
threshold is 2.0 days for Landfill bins, 1.8 days for Compostables, and 2.0 days for
Recyclables.

The urgency label system is intentionally simpler than the underlying model outputs. This
design reflects feedback from the project advisor that the tool must be usable by
operations staff without a technical background. The six-level scale mirrors the existing
red/green indicator in CLEAN and is intended to augment, not replace, the current driver
workflow.

\subsection{Short-Term Operational Recommendations}
Several immediate operational changes can be implemented to capture value from the forecasting system without requiring major changes to existing infrastructure or workflows. These recommendations focus on reducing unnecessary collections, improving prioritization, and aligning daily operations with predicted demand.

First, the operations team can implement a daily “do not collect” list using the LOW and MONITOR urgency categories. These bins are unlikely to require service within the current week and should be explicitly excluded from routes unless operationally convenient. Given that 35.1\% of historical collections occurred before bins reached threshold , even partial adherence to this rule could significantly reduce wasted trips.

Second, route planning can be guided using the urgency labels as a lightweight prioritization system. Bins classified as OVERDUE and HIGH should be treated as mandatory stops, while ELEVATED bins should form the core of near-term route planning. MODERATE bins can be included opportunistically when drivers are already nearby. This approach allows for improved efficiency without requiring full route optimization or new software systems.

Third, weekly planning can be introduced using the model’s strong performance over a seven-day horizon. Because the forecast correctly identifies approximately 95\% of bins that will require service within the next week , operators can generate a weekly service plan that distributes workload more evenly across days. This reduces reliance on reactive dispatching and helps prevent both overflow and over-servicing.

Fourth, a simple short-horizon validation rule should be used to address known limitations of the model at the 24-hour level. For bins with unusually long elapsed time since last collection relative to their historical average, operators should cross-check the CLEAN dashboard before deferring service. This hybrid approach combines predictive insight with real-time data to mitigate uncertainty.

Fifth, a small set of high-variability bins—particularly within the recyclables stream—should be monitored more closely. These bins exhibit higher prediction error and are more likely to deviate from expected patterns. Maintaining a manual watchlist for these locations provides a low-cost way to improve service reliability.

Finally, the organization should pilot a modest increase in the collection threshold for selected bins. External benchmarking indicates that similar systems operate effectively at higher thresholds (e.g., 75–80\%), suggesting that the current 60\% threshold may be conservative. A controlled pilot on low-risk bins would allow the team to evaluate the trade-off between efficiency gains and overflow risk.

\subsection{Urgency Band Distribution}

Applying the calibrated operational forecast to the most recent data timestamp produces a concrete prioritization across the 248-bin fleet. Of the 248 bins, 53 are classified as OVERDUE, 18 as HIGH, 42 as ELEVATED, 14 as MODERATE, 25 as LOW, and 96 as MONITOR. The OVERDUE and HIGH categories identify bins that should be collected immediately or within the next 24 hours. ELEVATED and MODERATE identify bins that are likely to require service within the near-term planning horizon and therefore form the core of current route planning.

The large MONITOR category does not indicate model failure. Instead, it reflects the structure of the fleet and the limits of the available data. Some bins lack a usable fitted curve, while others have fitted curves whose asymptotes remain below the 60\% collection threshold under normal conditions. For these bins, the forecast appropriately recommends periodic review rather than routine service. Operationally, the urgency-band system therefore separates the fleet into three groups: bins requiring immediate action, bins requiring near-term scheduling attention, and bins that can safely be monitored or deprioritized.

\begin{table}[ht]
\centering
\caption{Urgency band distribution at the most recent data timestamp.}
\label{tab:urgency_band_distribution}
\begin{tabular}{l c p{8cm}}
\hline
\textbf{Urgency Label} & \textbf{Count} & \textbf{Operational Meaning} \\
\hline
OVERDUE  & 53 & Already at or above threshold; collect immediately \\
HIGH     & 18 & Optimistic forecast reaches threshold within 24 hours \\
ELEVATED & 42 & Expected to reach threshold within 3 days \\
MODERATE & 14 & Expected to reach threshold within 5 days \\
LOW      & 25 & No collection needed within the current week \\
MONITOR  & 96 & No usable curve fit or asymptote below threshold; periodic review only \\
\hline
\textbf{Total} & \textbf{248} & Entire fleet \\
\hline
\end{tabular}
\end{table}

\subsection{Limitations}

The coarse 20\% fill-level measurement is the largest constraint on model accuracy. Fill
levels are reported only in steps of 0, 20, 40, 60, 80, and 100\%, introducing noise that
cannot be reduced through modeling. This is reflected in the average goodness-of-fit of
0.23 for the fill-curve fits.

The short-horizon limitation in Table~\ref{tab:horizons} is a direct consequence of the
same data constraint - without fill-level readings between collection events, the model
has no signal about intra-day fill dynamics or sudden surges. Weather features contributed
negligibly to prediction accuracy, reflecting Berkeley's mild and consistent climate.
Campus event features had modest effects on aggregate demand prediction but their impact
at the individual bin level is difficult to isolate.

The P10-P90 uncertainty bounds produced by the quantile model initially showed 59\%
empirical coverage against an 80\% target, indicating the raw intervals were too narrow.
This was corrected using split conformal calibration: a held-out calibration set (2024 H2)
was used to compute a symmetric correction factor of 17 hours, which is applied to all
intervals before operational use. After correction, empirical coverage on the 2025 test
set is 80.8\%, meeting the target. Corrected intervals have a median width of 5.4 days.
The bounds can be interpreted as approximate 80\% prediction intervals - for a given bin,
the true next collection time falls within the stated range approximately 80\% of the
time. Coverage is lower at short horizons (24-hour coverage remains limited for the same
structural reasons that constrain short-horizon recall), so the bounds are most reliable
for the 72-hour to one-week planning window.

The 14 low-usage bins without fill-curve fits are flagged as MONITOR and should be
incorporated into the fill-curve dataset as sufficient observation history accumulates.
Operational changes - such as the reusable container program at MLK Student Union - may
alter fill patterns for affected bins, which would require retraining on post-change data.

\newpage

% ── SUMMARY AND CONCLUSIONS ───────────────────────────────────────────────────
\section{Summary and Conclusions}

This project developed a forecasting system to help Cal Zero Waste shift from reactive
to proactive bin collection scheduling. Using 59,111 collection events from January 2023
through February 2026, the team built and evaluated multiple models, combined their
outputs into a rolling daily forecast, and validated the results against held-out 2025
test data.

The analysis confirmed that 35.1\% of all collections are wasted trips - bins emptied
well before reaching the 60\% collection threshold - and that the average bin is only
47.1\% full at the time of collection. The wasted trip decomposition shows that
approximately 77\% of these trips are potentially eliminable through better forecasting
and scheduling. The unexplained wasted trip rate has grown from approximately 8\% of all
collections in 2023 Q3 to approximately 29\% in early 2026, indicating that the gap
between current scheduling practices and actual bin needs has widened over time and
warrants further investigation.

The forecasting system performs well at the weekly planning horizon. The quantile
regression model reduces prediction error by 16\% relative to the naive baseline (mean
error of 40.3 hours versus 48.3 hours) and correctly identifies 95\% of bins that will
need collection within the next seven days. Short-horizon prediction remains limited by
the coarse 20\% fill-level measurement: at 24 hours, the model correctly identifies only
12\% of bins that will need service the following day. The fill-curve model addresses
this gap by estimating each bin's current fill level from its historical trajectory and
is used for near-term urgency classification in the 0 to 48 hour window.

The binary readiness classifier correctly identifies 85\% of bins that currently need
collection when tuned for high recall, with an overall discrimination score of 0.857.
The daily demand forecast provides fleet-level workload estimates for staffing and truck
scheduling, with mean errors of 6.8 to 7.7 collections per day across waste streams.

The uncertainty bounds accompanying each forecast have been calibrated to 80.8\%
empirical coverage using a held-out calibration set, with a median interval width of
5.4 days. The bounds are most reliable for the 72-hour to one-week planning window.

The operational output combines these components into a six-level urgency system and a
service frequency classification. Applied to the most recent data, the forecast flags
53 bins as OVERDUE and 18 as HIGH - requiring immediate or same-day service - and
classifies 25\% of the network as multi-visit bins requiring more than one collection
per week.

\subsection{Recommendations}

The following recommendations are organized by implementation timeline.

In the near term, Cal Zero Waste can reduce wasted trips without new infrastructure by
using the LOW and MONITOR urgency labels to generate a daily exclusion list - bins that
should not be visited unless operationally convenient. OVERDUE and HIGH bins should be
treated as mandatory route stops; ELEVATED bins should form the core of near-term
planning; MODERATE bins can be included opportunistically. A weekly service plan
generated from the seven-day forecast would distribute workload more evenly and reduce
reliance on reactive dispatching. For bins with unusually long elapsed time relative to
their historical average, operators should cross-check the CLEAN dashboard before
deferring service. A pilot program to modestly raise the collection threshold - for
example, from 60\% to 70\% on selected low-risk bins - is also worth evaluating, as
similar systems operate effectively at higher thresholds.

Over a longer horizon, models should be retrained quarterly as new CLEAN export data
accumulates. The growing unexplained wasted trip rate warrants investigation: it may
reflect unrecorded operational changes, route adjustments, or a behavioral shift in how
drivers use the CLEAN dashboard. As Project 10B matures, integrating the urgency labels
and calibrated forecast bounds into a route optimization system is the most direct path
to achieving the division's 90\% collection efficiency target. The 14 bins currently
flagged as MONITOR should be incorporated into the fill-curve dataset as observation
history grows.

\newpage

% ── SUSTAINABLE IMPLEMENTATION ────────────────────────────────────────────────
\section{Sustainable Implementation}

Keeping the forecasting system operational requires three ongoing activities: periodic
model retraining, structured data export, and procedures for handling changes to the
bin network.

\subsection{Retraining Cadence}

The quantile regression and binary classifier models should be retrained quarterly
using accumulated CLEAN export data. Fill-curve coefficients are more stable and can
be updated semi-annually, or immediately when a bin's usage pattern changes materially.
Retraining requires exporting the updated collection log, running the data processing
and modeling notebooks in sequence, and replacing the output files used by the
operational forecast. The full pipeline can be completed in approximately two to three
hours on a standard laptop.

\subsection{CLEAN Export Process}

The primary data source is the CLEAN Management Console export, currently pulled by
the project sponsor on an as-needed basis. For the forecast to remain accurate, exports
should occur on a fixed schedule - weekly or bi-weekly - and be appended to the master
collection log before running the forecast update. The export format and column mapping
are documented in the data processing notebook.

\subsection{Staff Interface}

The operational forecast output is designed for use by operations staff without a
technical background. The urgency labels and frequency classifications can be consumed
directly from the forecast output file, or incorporated into the existing CLEAN
dashboard workflow as a supplementary priority list. No model interpretation is required
at the point of use.

\subsection{Handling Operational Changes}

Significant changes - such as the reusable container program introduced at MLK Student
Union, bin relocations, or changes to the collection threshold - can shift fill patterns
for affected bins. When such changes occur, the affected bins should be flagged and their
fill-curve coefficients reset once at least two months of post-change collection data
are available. The quantile regression and binary classifier models will self-correct
through quarterly retraining, but predictions for affected bins should be treated with
additional uncertainty in the interim.

\newpage

% ── INTEGRATION WITH PROJECT 10B ──────────────────────────────────────────────
\section{Integration with Project 10B}

Project 10B is a companion effort focused on route dispatch optimization. It will
consume the forecasting outputs from this project to generate optimized daily collection
routes. Coordination with Team 10B is ongoing; the interface specification below
reflects the current state of planning and will be updated once that coordination is
complete.

The primary handoff artifact is the operational forecast output file, updated on a
rolling daily basis. The fields most relevant to dispatch optimization are the urgency
label, the calibrated optimistic, median, and conservative time-to-collection estimates,
the estimated current fill level, the service frequency classification, and the bin
location and waste stream identifiers needed for route assignment.

For planning purposes, Team 10B can treat the urgency labels as a soft constraint on
route construction: OVERDUE and HIGH bins are mandatory daily stops; ELEVATED and
MODERATE bins are candidates for inclusion based on route capacity; LOW and MONITOR bins
are optional stops. The calibrated time-to-collection estimates provide additional
flexibility - a bin with a conservative estimate of 36 hours can be safely deferred to
the following day, while a bin with an optimistic estimate of 12 hours likely cannot.

The update cadence, file format, and exact field mapping will be documented in a formal
interface specification once Team 10B's routing model requirements are confirmed.

\newpage

% ── REFERENCES ────────────────────────────────────────────────────────────────
\section*{References}
\addcontentsline{toc}{section}{References}

\begin{itemize}[label={}, leftmargin=0pt, itemindent=0pt, itemsep=6pt, parsep=0pt]

  \item Hasan, M., et al. (2024). Prophet time series modeling of waste disposal rates in
  four North American cities. \textit{Environmental Science and Pollution Research},
  Springer.

  \item Taylor, S.J., \& Letham, B. (2018). Forecasting at scale.
  \textit{The American Statistician}, 72(1), 37-45.

  \item Chen, T., \& Guestrin, C. (2016). XGBoost: A scalable tree boosting system.
  \textit{Proceedings of the 22nd ACM SIGKDD International Conference on Knowledge
  Discovery and Data Mining}, 785-794.

  \item Ke, G., et al. (2017). LightGBM: A highly efficient gradient boosting decision
  tree. \textit{Advances in Neural Information Processing Systems}, 30.

\end{itemize}

\newpage

% ── APPENDIX A: TECHNICAL DETAILS ─────────────────────────────────────────────
\appendix
\section{Technical Details}

This appendix contains the mathematical details and model specifications referenced in the
main body. It is intended for readers with a quantitative background.

\subsection*{Fill-Curve Equations}

For each bin, the best-fitting curve from the following set is selected based on $R^2$:

\begin{itemize}
  \item \textbf{Sigmoid:} $F(t) = \dfrac{L}{1 + e^{-k(t - t_0)}}$,
  where $L$ is the asymptote, $k$ is the growth rate, and $t_0$ is the inflection point.
  Time to threshold: $T^* = t_0 - \dfrac{\ln(L/60 - 1)}{k}$

  \item \textbf{Saturating exponential:} $F(t) = a(1 - e^{-bt})$

  \item \textbf{Power law:} $F(t) = a \cdot t^b$

  \item \textbf{Logarithmic:} $F(t) = a \ln(t) + b$

  \item \textbf{Linear:} $F(t) = a \cdot t$
\end{itemize}

In all cases, $t$ is hours elapsed since the last collection event and $F(t)$ is predicted
fill level as a percentage of capacity. Coefficients are fitted per-bin using nonlinear
least squares. Mean $R^2$ across all 224 fitted bins is 0.23, reflecting the 20\%
measurement quantization rather than model inadequacy.

\subsection*{Machine Learning Model Details}

\textbf{Train/test split:} 2023-2024 for training, 2025 onward for testing.

\textbf{Features:} Day of week, month, cyclical time encodings, hours elapsed since last
collection, last observed fill level, 7-day and 28-day rolling averages of elapsed time
and fill level, campus event flags (semester, spring recess, high-impact events, move-in
week), weather (daily max/min temperature, precipitation, wind speed), location encoding,
stream type, and fill-curve statistics from Phase 1 (median elapsed hours, predicted fill
at 24h, 48h, 72h).

\textbf{XGBoost binary classifier:} Class imbalance addressed with
\texttt{scale\_pos\_weight} proportional to class ratio. Classification threshold selected
at the lowest value yielding recall $\geq 0.85$ on the precision-recall curve (threshold
= 0.33). Early stopping with 3-fold time-series cross-validation.

\textbf{LightGBM quantile regression:} Separate models trained at $\alpha \in
\{0.10, 0.50, 0.90\}$ using the pinball loss. 600 estimators, learning rate 0.05,
63 leaves, feature and bagging fraction 0.8.

\textbf{Prophet:} Weekly and yearly seasonality. Exogenous regressors:
\texttt{is\_weekend}, \texttt{is\_semester}, \texttt{is\_high\_impact\_event},
\texttt{is\_spring\_recess}. Fit separately per stream type on daily aggregate needed
collection counts.

\end{document}
